\section{Related Work}
\label{submission:related}

\subsection{Weight-Space Model Merging}
Model merging in weight space represents a highly efficient paradigm for multi-task learning and model editing without additional training costs \cite{wortsman2022model, mcmahan2017communication}. Early works established that models sharing the same pre-trained initialization lie within the same low-loss basin and can be linearly interpolated \cite{mirzadeh2020linear, gupta2020stochastic}. Task Arithmetic \cite{ilharco2023editing, chane2022patching} formalizes this by extracting \emph{task vectors}---defined as the difference between the fine-tuned expert weights and the pre-trained base weights---and adding them to the base model. To improve linear blending, Fisher-weighted averaging \cite{matena2022merging} and regression-based weight alignment \cite{jin2023regmean} utilize local quadratic approximations or training data statistics to weight parameter updates. While Task Arithmetic works well for similar tasks, linear addition causes catastrophic representational collapse when task vectors correspond to highly conflicting or orthogonal domains, as updates in opposing directions cancel each other out.

\subsection{Interference Mitigation in Model Merging}
To resolve weight-space interference, several filtering and routing heuristics have been proposed. TIES-Merging \cite{yadav2023ties} trims small-magnitude updates, resolves sign conflict by electing a majority sign direction at each coordinate, and averages the remaining updates that agree with the elected sign. DARE \cite{yu2024dare} utilizes random dropout to sparse-out task updates, followed by uniform scaling to preserve the original representation's scale, before performing standard averaging. Other approaches attempt to align model parameters via permutation symmetries \cite{ainsworth2022git}, resolve representation mismatch after merging \cite{jordan2022repair}, or merge models from disparate tasks by aligning activation distributions \cite{stoica2023zip}. 
While TIES-Merging and DARE mitigate sign conflicts, they still ultimately \emph{average} or blend parameter updates across models. Under severe task conflict, this weight blending continues to cause representation dilution and performance collapse. In contrast, Exclusive Parameter Merging (EPM) limits weight-space interference by routing updates based on coordinate-wise exclusivity: each parameter update is primarily allocated to its dominant expert, while a controlled background blend retains general representation alignment. This mitigates catastrophic signal erasure without the heavy optimization overhead of continuous layer-wise scaling frameworks.

\subsection{Pruning and Sparsity}
Weight-space pruning has a rich history as a regularization technique to compress models and reduce overfitting \cite{han2015learning, frankle2018lottery}. In model merging, pruning can be applied as a pre-processing step (Prune-then-Merge) or co-optimized during merging. ZipMerge co-optimizes layer-wise magnitude-pruning thresholds and continuous scaling coefficients during test-time. However, joint layer-wise search is computationally heavy and highly unstable. EPM approaches sparsity with extreme simplicity: we apply a standard global threshold on the exclusive task vectors to achieve overall target sparsities (e.g., 50\% or 80\%) without needing any layer-wise tuning, showing that structured coordinate exclusivity provides immense resilience to heavy pruning.

\subsection{The Overfitting-Optimizer Paradox in Test-Time Tuning}
A major trend in modern model merging is the test-time adaptation (TTA) of continuous merging coefficients using unsupervised loss functions (like Shannon entropy minimization) on small calibration batches. While this continuous optimization minimizes the calibration objective, it suffers from the \textbf{Overfitting-Optimizer Paradox}. As analyzed in PolyMerge and OFS-Tune, high-dimensional layer-wise tuning (often featuring hundreds of continuous coefficients) transductively overfits to the small calibration sample, destroying the generalizable feature representations of the model and degrading multi-task test accuracy. Alternative search approaches like evolutionary algorithms \cite{akiba2024evolutionary} optimize merging recipes globally, though they still require substantial test-time validation sets or search time. EPM avoids this completely by restricting the optimization space to only $K$ global task-level coefficients, using a zero-order (1+1) Evolution Strategy on a tiny offline validation split. This minimal parameterization maintains perfect stability and ensures high generalizability.
